\section*{Report on Use of AI Tools}
OpenAI Codex was used during November 2026 to assist with translation, model formulation, Python implementation, source research, English drafting, LaTeX preparation and figure generation. The user also supplied a Google Gemini analysis as a PDF excerpt; its process-oriented explanation informed the early organization. Gemini was not used as a scientific source for the physical equations or numerical inputs.

\noindent\textbf{Development record.} The assistance covered freezer service accounting; psychrometric and SET calculations; the shared local-delivery operator; station heat, moisture and surface balances; candidate deployment comparisons; scope separation between efficiency and effectiveness; and the final CP-SET restructuring. The supplied drawings and earlier award papers were used for context and presentation, with their data provenance distinguished from scenario inputs.

\noindent\textbf{Latest revision.} The user selected Fremont Street Experience and the Honolulu Zoo ticketing queue for a large--small venue comparison. Codex researched site identity and climate, specified the assumption table, implemented the shared CP-SET calculation, and generated climate-transfer, scale and boundary tests. The revision uses an analytical efficiency bound and integer cost optimization, with a separate enumeration audit for the small venue. Associated summaries, operating guidance and standalone PDF figures were updated from the computed results. The Python implementation, input files, CSV outputs, figure scripts and verification records are retained with the LaTeX project.

\noindent\textbf{Verification record.} Automated checks cover heat and water conservation, service-energy identities, limiting-case recovery, time-step convergence and consistency between tables, figures and outputs. The AI-assisted workflow performed these checks; independent team review is recorded separately by the team.

\noindent\textbf{Final editorial revision.} Codex expanded the introduction, added model evaluation, normalized decimal notation, revised the public blog, placed updated user-supplied architectural drawings and adjusted figure terminology. Core numerical outputs and optimization results were retained. All AI-tool descriptions and references in this manuscript are consolidated in this report at the user's request.

\noindent\textbf{Visualization revision.} Codex wrote reproducible Python scripts for continuous station-field reconstruction, a NASA POWER climate-context map, and separate energy/service accounts with venue radar comparisons. Assumed spatial kernels and coordinates are disclosed; downloaded climate metadata are archived. The existing model inputs, numerical outputs and optimization results were frozen and hash-checked. Additional plots were consolidated and the compiled PDF was inspected for layout.

\noindent\textbf{Tool identification and transcript archive.} OpenAI Codex, accessed November 8, 2026; the session identifies Codex as based on GPT-6; a more specific deployment identifier was not supplied. Google Gemini, user-supplied conversation excerpt; the model identifier is not present in that excerpt. The original prompts and complete retained outputs are held in the associated conversation records. This narrative is an activity record, not a replacement for those transcripts.

\noindent\textbf{Policy reference.} COMAP, \emph{Use of AI Tools in COMAP Contests}, v102025, \url{https://www.contest.comap.com/undergraduate/contests/mcm/flyer/Contest_AI_Policy.pdf}; and 2026 HiMCM instructions, \url{https://www.contest.comap.org/highschool/contests/himcm/instructions.html}, accessed November 8, 2026.
